% ===== 原文 pp.20–21 =====
% tr03a：pp.20–21。注意：p.21 末三行（定理 9 的 "Proof." 开头及其短正合列）并入 tr03b，
% 以免 amsthm 的 proof 环境被 chunk 文件边界切断（check.py 逐文件核对 \begin/\end 配对）。
% p.19 末句由 tr-02 写完，本片从 p.20 首 "Write the length identity as" 起。
% [原文 p.20]

把长度恒等式写成
\[
L=\sum_{\alpha}\cD(L,\ell_\alpha,\ell_\alpha)=\sum_{h\in\MCG/\Stab(\alpha)}\cD(L,\ell_{h(\alpha)},\ell_{h(\alpha)}),
\]
利用 MCG 对测地长度函数的作用，得到
\[
LV(L)=\int_{\cT(L)/\MCG}\sum_{\MCG/\Stab(\alpha)}\cD(L,\ell_\alpha\circ h^{-1},\ell_\alpha\circ h^{-1})\,\omega,
\]
在 $\cT$ 上作变量替换 $p=h(q)$，得到
\[
\sum_{h\in\MCG/\Stab(\alpha)}\int_{h(\cT(L)/\MCG)}\cD(L,\ell_\alpha,\ell_\alpha)\,\dd\tau\,\dd\ell
=\int_{\cT(L)/\Stab(\alpha)}\cD(L,\ell_\alpha,\ell_\alpha)\,\dd\tau\,\dd\ell,
\]
再利用 $\Stab(\alpha)$ 的基本域，得到积分
\[
\int_0^\infty\int_0^\ell\cD(L,\ell,\ell)\,\dd\tau\,\dd\ell.
\]
对 $\tau$ 的积分给出一个因子 $\ell$。

导数 $\partial\cD(x,y,z)/\partial x$ 与 $\partial\cR(x,y,z)/\partial x$ 比原来的函数 $\cD$ 与 $\cR$ 更简单——利用这一观察并对 $L$ 求导，即得到 $LV(L)$ 的导数公式
\[
\frac{\partial}{\partial L}LV(L)=\int_0^\infty\frac{1}{1+e^{\ell+\frac{L}{2}}}
+\frac{1}{1+e^{\ell-\frac{L}{2}}}\;\ell\,\dd\ell=\frac{\pi^2}{6}+\frac{L^2}{8}.
\]
于是得到公式 $V(L)=\frac{\pi^2}{6}+\frac{L^2}{24}$。

% [Tutor Note: 上式之所以能由"对 L 求导"一步算出，是因为 p.18 公式 (3) 给出了
%  $\partial\cD(x,y,z)/\partial x=H(y+z,x)$：取 $x=L$、$y=z=\ell$，得
%  $\partial\cD(L,\ell,\ell)/\partial L=H(2\ell,L)=\frac{1}{1+e^{\ell+L/2}}+\frac{1}{1+e^{\ell-L/2}}$，
%  正是本式的被积函数。可核对的数值锚点：$L=1$ 时右端 $\int_0^\infty[\cdots]\ell\,\dd\ell=1.76993407$，
%  而 $\pi^2/6+1/8=1.76993407$，完全吻合；$L=0$ 时积分化为
%  $2\int_0^\infty\ell/(1+e^{\ell})\,\dd\ell=2\cdot\pi^2/12=\pi^2/6$，与 $V(0)=\pi^2/6$ 自洽。
%  再由 $\partial_L(LV)=\pi^2/6+L^2/8$ 积分回去并用 $V$ 在 $L=0$ 的已知值定常数，即得
%  $LV=\pi^2L/6+L^3/24$，故 $V=\pi^2/6+L^2/24$。]

我们为一般情形作准备。在代数几何的相交理论中，椭圆对合（elliptic involution）会导致椭圆相交计数要乘以因子 $1/2$。这个因子对应于：万有椭圆曲线（universal elliptic curve）的通用纤维是椭圆曲线对其对合的商。沿着这条思路，若把 $V_{1,1}(L)$ 形式地定义为给定值 $V(L)$ 的 $1/2$，则一般的体积递推会得到简化。在映射类群理论中，对界出一个环面的简单闭曲线，椭圆对合表现为\emph{半 Dehn 扭转}（half Dehn twist）。具体地说，考虑曲面的基本群 $\pi_1(R)$，采用标准表现 $a_1b_1a_1^{-1}b_1^{-1}\cdots a_gb_ga_g^{-1}b_g^{-1}c_1\cdots c_n=1$，其中 $c_j$ 是绕第 $j$ 条边界的回路。关于曲线 $a_1b_1a_1^{-1}b_1^{-1}$ 的半 Dehn 扭转是 $\pi_1(R)$ 的如下自同构：$a_1\mapsto b_1a_1^{-1}b_1^{-1}$，$b_1\mapsto b_1^{-1}$；$a_j\mapsto a_1^{-1}a_ja_1$，$b_j\mapsto a_1^{-1}b_ja_1$，对 $j=2\ldots g$；以及 $c_j\mapsto a_1^{-1}c_ja_1$，对 $j=1\ldots n$。半 Dehn 扭转的平方是 Dehn 扭转，而半 Dehn 扭转在相应的 FN 参数上的作用为 $(\ell,\tau)\mapsto(\ell,\tau+\ell/2)$。

现在为余体积公式作准备。设 $R$ 是带测地边界 $\beta_1,\ldots,\beta_n$ 的双曲曲面。考虑一条加权多重曲线（weighted multicurve）
\[
\gamma=\sum_{j=1}^m a_j\gamma_j,
\]
% [原文 p.21]
其中 $a_j$ 为实权重，$\gamma_j$ 是互不相同、彼此不交的简单闭测地线。定义 $\Stab(\gamma)\subset\MCG$ 为稳定这族无标号、带权重的测地线的映射类所成之子群——$\Stab(\gamma)$ 的元素可以置换多重曲线中权重相同的分支。把单条测地线的稳定化子记为 $\Stab(\gamma_j)$，把其中保持定向的元素所成之子群记为 $\Stab_0(\gamma_j)$。

把沿 $\gamma$ 剪开所得的曲面记为 $R(\gamma)$——每条 $\gamma_j$ 产生两条新边界——$R(\gamma)$ 可以是断开的。剪开后的曲面，其（乘积）Teichmüller 空间记为 $\cT(R(\gamma);\mathbf{x})$，$\mathbf{x}=(x_1,\ldots,x_m)$，其中 $\gamma_j$ 对应的一对边界具有长度 $x_j$。把 $R(\gamma)$ 各分支的映射类群之积记为 $\MCG(R(\gamma))$，相应的模空间之积记为 $\cT(R(\gamma);\mathbf{x})/\MCG(R(\gamma))$。对 $\cT(R(\gamma);\mathbf{x})$ 上对应于 $R(\gamma)$ 各分支的辛形式之积，其体积 $V(R(\gamma);\mathbf{x})$ 等于各分支模空间体积之积，这里同样要求 $\gamma_j$ 对应的一对边界具有公共长度 $x_j$。讨论中还涉及有限对称群
\[
\operatorname{Sym}(\gamma)=\Stab(\gamma)/\bigcap_j\Stab_0(\gamma_j),
\]
即可能置换 $\gamma$ 各元素、并反转其定向的映射类所成之群。

把一个函数在一个群上的各平移求和，就得到群作用下不变的函数。从一个在无穷远处适当小的函数 $f$ 出发，引入 MCG 和
\[
f_\gamma(R)=\sum_{\MCG/\Stab(\gamma)}f(\sum_{j=1}^{m}a_j\ell_{h(\gamma_j)}(R))
\eqno{(4)}
\]
下一个定理把模空间积分
\[
\int_{\cM(R)}f_\gamma\,\dd V
\]
表示为低维模空间体积的加权积分。

\begin{thm}{9} [Mir07a, Thrm. 7.1] 余体积公式（covolume formula）。设 $\gamma=\sum_{j=1}^m a_j\gamma_j$ 为加权多重曲线，$f$ 为在无穷远处适当小的函数，$f_\gamma$ 为 $f$ 的 MCG 和，则
\[
\int_{\cT(R)/\MCG}f_\gamma\,\dd V=(|\operatorname{Sym}(\gamma)|)^{-1}\int_{\R^m_{>0}}f(|\mathbf{x}|)V(R(\gamma);\mathbf{x})\,\mathbf{x}\cdot\dd\mathbf{x},
\]
其中 $|\mathbf{x}|=\sum_j a_jx_j$，$\mathbf{x}\cdot\dd\mathbf{x}=x_1\cdots x_m\dd x_1\cdots\dd x_m$。
\end{thm}

% [Tutor Note: 定理 9 有两处最易漏：因子 $(|\operatorname{Sym}(\gamma)|)^{-1}$ 与测度
%  $\mathbf{x}\cdot\dd\mathbf{x}$。(i) 当 $m=1$、曲面为亏格 1 单边界、$\gamma=a_1\gamma_1$ 时，
%  p.19 已说明 $\Stab(\gamma_1)$ 是 Dehn 扭转与对合 $\Z/2\Z$ 的半直积，故
%  $|\operatorname{Sym}(\gamma_1)|=2$，右端带 $1/2$——这正是 p.22 末"对单边界环面改用 $1/2$ 的
%  $V(L)$"的同一因子；p.32 记 $V_{1,1}(L)=\pi^2/12+L^2/48$ 并注明"已含椭圆对合 $1/2$ 因子"，
%  而 p.20 算出的 $\pi^2/6+L^2/24$ 是未含该因子的值（恰为其 2 倍），可逐项对照。
%  (ii) $\mathbf{x}\cdot\dd\mathbf{x}=x_1\cdots x_m\dd x_1\cdots\dd x_m$ 中的 $\prod_j x_j$
%  来自 p.22 证明里沿 $\tau_j$ 方向的积分：$\operatorname{Dehn}_*(\gamma_j)$ 只平移 $\tau_j$，
%  其基本域长度为 $\ell_j$（单边界环面的半扭转情形为 $\ell_j/2$，该 $1/2$ 另行计入体积值），
%  故每个 $j$ 贡献一个因子 $x_j$，合起来就是 $\prod_j x_j$；若只写
%  $\dd x_1\cdots\dd x_m$ 会整体差一个 $\prod_j x_j$，做递推时会对不上体积多项式。]
